\section{Naive Approaches}


%\subsection{Hamiltonian Neural Networks (HNN)}
\begin{frame}{Hamiltonian Neural Networks (HNN)}
    Choosing the Hamiltonian would be the \textbf{best option} but a set of \textbf{canonical variables} is needed.
    
    \begin{block}{Canonical Variables}
        \begin{columns}[c] % [c] allinea verticalmente al centro
            
            % Colonna di sinistra: Equazione
            \begin{column}{0.45\textwidth}
                \[
                    \mathbf{z} = 
                    \begin{bmatrix}
                        \mathbf{q} \\
                        \mathbf{p}
                    \end{bmatrix}
                \]
            \end{column}

            % Colonna di destra: Legenda
            \begin{column}{0.50\textwidth}
                \begin{itemize}
                    \item $\mathbf{q}$: generalized displacements
                    \item $\mathbf{p}$: generalized momenta
                \end{itemize}
            \end{column}
            
        \end{columns}
    \end{block}

    Given the \textbf{Hamiltonian} as $\mathcal{H}$, it's possible to find the dynamics with: 

    $$\frac{dq}{dt} = \frac{d\mathcal{H}}{dp} \hspace{2em} \frac{dp}{dt} = -\frac{d\mathcal{H}}{dq}$$
    
    It's possible to use an MLP neural network that learns the Hamiltonian $\mathcal{H}_\theta$ and a loss that enforces the dynamics: 

    $$\mathcal{L}_{HNN} =  \left \lVert \frac{d\mathcal{H}_\theta}{dp} - \frac{\partial q}{\partial t} \right \rVert_2 + \left \lVert \frac{d\mathcal{H}_\theta}{dq} + \frac{\partial p}{\partial t} \right   \rVert_2$$

    Two problems of momenta: 

    \begin{itemize}
        \item they depend on the parameters of the system
        \item more importantly they depend on the coupled dynamics not known
    \end{itemize}
    
    
\end{frame}


%\subsection{Parameters Predictor}
\begin{frame}{Entries Prediction}

    Remembering equation (4), is it necessary to know matrices structure thoroughly?

    \begin{block}{Matrices Structure}
        \par\textbf{Goal} $\rightarrow$  \textbf{Graybox} approach $+$ insights on system properties. 
        \par\textbf{GT (Ground Truth)} $\rightarrow$ \textbf{Newton-Euler} algorithms. \\
        \par\textbf{Application}: \textbf{Quick inference} for control purposes. \\
    \end{block}

    \begin{figure}
        \centering
        \includegraphics[width=0.7\textwidth]{06_naive_approaches/entries_pred.pdf}
        \caption{Entries prediction architecture}
    \end{figure}

\end{frame}

\begin{frame}{Entries Prediction - Considerations}
    \begin{itemize} 
        \item \textbf{Advantages}: Quick training, very good results.
        \item \textbf{Disadvantages}: Requires GT values for matrices entries.
        \item We would like to obtain the same results \textbf{without} knowing the entries.
    \end{itemize}

    \begin{figure}
        \centering
        \includegraphics[width=0.8\textwidth]{06_naive_approaches/v1000wnEntries/plots/evolution_comparison_1.png.png}
    \end{figure}

\end{frame}

\begin{frame}
    \begin{figure}
        \centering
        \includegraphics[width=1.0\textwidth]{06_naive_approaches/v1000wnEntries/plots/mat_comp_1.png.png}
    \end{figure}
\end{frame}





